\paragraph{Deriving a centered clock and constructing its new ground.}
The CPT obstruction motivates a separate, explicitly changed charge
clock. On a bipartite Weyl lattice with $V$ sites, the one-particle
antiunitary $J=\eta(i\sigma_y)\mathcal K$, where
$\eta_x=(-1)^{x_1+x_2+x_3}$ and $\mathcal K$ is conjugation, anticommutes
with the hopping. Its canonical particle-hole transform satisfies
\[
 \mathcal P(N)=2V-N,\qquad
 \mathcal P(H)=H+U(V-N),\qquad
 \mathcal P(H-\kappa N)-(H-\kappa N)=(U-2\kappa)(V-N).
\]
Thus $K=H-UN/2$ is the unique scalar charge centering invariant on the
full Fock algebra. This changes charged dynamics; the original Hamiltonian
and its earlier ground are not silently replaced.

The actual cube gives a rigorous changed-reference test. Its retained
all-number theorem states $H\ge-\sqrt3/4$. The old $N=4$ ground has
centered energy $-2-\sqrt3/4>-4$, while a singly occupied $N=8$ product
state has expectation $-4$. Moreover, for $N\le7$,
$K\ge-7/2-\sqrt3/4>-4$, and particle-hole symmetry transfers this
bound to $N\ge9$. Every new global centered ground therefore has
exactly eight particles. A product of four exact dimer grounds on a
perfect matching improves the variational bound: inter-dimer hopping
expectations vanish because each factor has fixed charge. Hence
\[
 -4-\frac{\sqrt3}{4}\le E_K\le-2-\frac{\sqrt{17}}2,
 \qquad
 \inf\operatorname{spec}(K|_{N\ne8})-E_K
 \ge\frac{\sqrt{17}}2-\frac{\sqrt3}4-\frac32>\frac18.
\]
These are exact full-sector bounds, not relativistic rest masses.

We additionally import Lieb's repulsive Hubbard ground-spin theorem,
in the form of Theorem~1, printed p.~8, of his
\href{https://arxiv.org/pdf/cond-mat/9311033}{1993 review}; see also the
\href{https://doi.org/10.1103/PhysRevLett.62.1201}{1989 result} and its
\href{https://doi.org/10.1103/PhysRevLett.62.1927.5}{erratum}.
The actual hidden-frame hopping is real, spin independent and connected,
with a balanced bipartition of four plus four sites and strictly positive
onsite repulsion. These finite hypotheses are checked from the source;
the theorem allows mixed real hopping signs. Together with the separately
proved $N=8$ ground sector, it gives a unique global centered cube ground
with hidden spin zero. The review's centered interaction convention
differs from $K$ only by the constant $2$. No quantitative neutral gap,
physical spin dictionary or thermodynamic uniqueness follows here.

The actual cube also has a faithful full-CAR one-site state. Write
$K=T+D-4$ on $N=8$. The source identity $A^2=3I$ gives
$\|T\|=\sqrt3/2$. If $P_0$ projects onto the 256 singly occupied
configurations, every nonzero hopping leaves this subspace, so
$P_0TP_0=0$. The source's signed vertex flips are transitive. Set
$p=\omega(1-P_0)\le\omega(D)$ and
$\delta=\sqrt{17}/2-2>3/50$. The block Cauchy--Schwarz estimate gives
\[
 |\omega(T)|\le2\|T\|\sqrt{p(1-p)}+\|T\|p
 \le3\|T\|\sqrt{\omega(D)},\qquad
 \omega(T)+\omega(D)\le-\delta.
\]
Thus the common site doublon expectation $d=\omega(D)/8$ satisfies
$d\ge\delta^2/54>1/15000$, while
$d\le(\|T\|-\delta)/8<1/8$. Gauge, hidden-spin and particle-hole
invariance determine the complete one-site state and its modular operator:
\[
 \rho_x=\operatorname{diag}(d,\tfrac12-d,\tfrac12-d,d),\qquad
 K_x=-\log(\tfrac12-d)I+
 \log\!\left(\frac{\tfrac12-d}{d}\right)(N_x-1)^2.
\]
In particular, $\lambda_{\min}(\rho_x)>1/15000$ and
$\|K_x\|<\log15000$. For every same-site density $\sigma$, put
$\Delta=\sigma-\rho_x$ and $\rho_\epsilon=\rho_x+\epsilon\Delta$.
The full-site entropy first law and remainder are
\[
 \left.\frac{dS(\rho_\epsilon)}{d\epsilon}\right|_0
 =\operatorname{Tr}(\Delta K_x),\qquad
 0\le D(\rho_\epsilon\Vert\rho_x)
 \le\frac{7500\epsilon^2}{1-\epsilon}\|\Delta\|_{\rm HS}^2,
 \quad 0\le\epsilon<1.
\]
This finite result uses no fitted ground vector and covers noncommuting
variations. It is not an infinite-ground faithfulness theorem or an
identification of the modular operator with a geometric charge.

For a dimer with hopping $t>0$, let $s=\sqrt{U^2+16t^2}$. Its unique
full-Fock centered ground is a singlet/doublon superposition with energy
$-(U+s)/2$. Its one-site density is
\[
 \rho_{\rm site}=\operatorname{diag}
 (d,\tfrac12-d,\tfrac12-d,d),\qquad d=(1-U/s)/4>0.
\]
Thus this changed reference is faithful on the full local CAR algebra.
For a fixed local spin, both addition and removal have energies and weights
\[
 g_\pm=s/2\pm t,\qquad w_\pm=(1\mp4t/s)/4.
\]
Each charge branch has total weight $1/2$; zero weights are retained.
At the actual bond parameters $U=1,t=1/16$, the positive energies are
$(2\sqrt{17}\pm1)/16$, and the signed second moment is $65/256$.

Normalized full centered-ground projectors, with a tracial exterior,
are particle-hole, gauge and hidden-spin invariant. Translation averaging
and the existing boundary/volume and local-dynamics estimates give a
new infinite-volume ground with positive centered GNS energy and exactly
one particle per site in expectation. The old unshifted state is not reused. The exact
canonical composite $d=(2n_{\rm opposite}-1)c$ gives
\[
 [K,c(f)]=-c(hf)-(U/2)d(f),\qquad \|d(f)\|=\|f\|.
\]
Consequently the same all-word compactness construction supplies a centered
eight-valley continuum net. Its moment identities are
$M_1=h$, $M_2=h^2+U^2/4$; only the corresponding second-moment upper
bound is passed to the continuum. Particle-hole symmetry pairs the
complementary valleys and equates averaged addition/removal measures.
This resolves the earlier first-moment obstruction, without establishing
full spacetime CPT or Lorentz covariance.

There is also a stronger microscopic occupation bound. On a positive
kinetic band with $\epsilon\ge\delta>0$, the GNS maps
$Cf=c(P_+f)\Omega$, $Df=d(P_+f)\Omega$ obey
$LC+C\epsilon=-(U/2)D$. Integrating
$e^{-tL}Ce^{-t\epsilon}$ and using $L\ge0$, $\|D\|\le1$ gives
$\|C\|\le U/(2\delta)$. Translation invariance, arbitrary measurable
band windows and rational thresholds yield
$w(k)\le U^2/(2c^2\sum_j\sin^2k_j)$, including the negative-band
hole estimate. Folding into the node cells and integrating over a
circumscribed ball gives
\[
 \int\frac{d^3k}{(2\pi)^3\sum_j\sin^2k_j}
 \le\frac{\pi\sqrt3}{2}<\frac{11}{4},\qquad
 D\le\min\{2,11(U/c)^2/8\}.
\]
The same argument on the original uncentered ground uses
$\|B\Omega\|\le\sqrt r$ and $r\le1/2$, giving the separate bound
$D\le\min\{2,11(U/c)^2/4\}$. Neither bound is uniform on a normalized
packet in a shrinking momentum region. Native preparation of the new
ground, surviving continuum interactions, measured particle identities
and the common geometric Einstein premises remain further requirements.
The exact source checks and proofs are in
\texttt{model/centered\_source\_ground}.
